% Part: first-order-logic
% Chapter: sequent-calculus
% Section: proving-things-quant

\documentclass[../../../include/open-logic-section]{subfiles}

\begin{document}

\olfileid{fol}{seq}{prq}

\olsection{संख्यापकांसह \usetoken{P}{derivation}}

\begin{ex}
क्रमवर्ती $\lexists[x][\lnot !A(x)]
\Sequent \lnot \lforall[x][!A(x)]$ ची एक $\Log{LK}$-!!{derivation} द्या.

संख्यापक हाताळताना आयगेन चराच्या अटीचा भंग होणार नाही याची खात्री करावी
लागते; काही वेळा त्यासाठी ठरावीक अनुमाने कोणत्या क्रमाने करायची, यात फेरबदल
करावा लागतो. सामान्यतः आयगेन चराची अट लागू असलेले नियम आधी हाताळण्याचा
प्रयत्न केल्यास मदत होते (पूर्ण झालेल्या सिद्धतेत ते खालच्या बाजूस असतील).
तसेच पुढे काय येईल याचा विचार करून आरंभीची क्रमवर्ती कशी असेल याचा अंदाज
लावणे उपयुक्त ठरते. आपल्या बाबतीत ती $!A(a) \Sequent !A(a)$ यासारखी असावी
लागेल. याचा अर्थ असा की संख्यापकांचे नियम ``उलट्या दिशेने'' लागू करताना
$\lforall$ नियम आणि $\lexists$ नियम या दोन्हींसाठी एकच पद---ज्याला आपण $a$
म्हणू---निवडावे लागेल. प्रत्येक नियमासाठी वेगवेगळी पदे निवडली, तर शेवटी
$!A(a) \Sequent !A(b)$ यासारखी क्रमवर्ती मिळेल; ती अर्थातच निष्पन्न करता
येत नाही.

नेहमीप्रमाणे आपण पुढील क्रमवर्ती लिहून सुरुवात करतो:
\begin{prooftree}
\AxiomC{}
\UnaryInf$\lexists[x][\lnot !A(x)] \fCenter \lnot \lforall[x][!A(x)]$
\end{prooftree}
आपण \LeftR{\exists} नियम किंवा \RightR{\lnot} नियम यांपैकी कोणताही लागू करू
शकतो. \LeftR{\exists} नियमाला आयगेन चराची अट लागू असल्यामुळे तो उशिरा
हाताळण्यापेक्षा लवकर हाताळणे योग्य ठरेल; म्हणून तोच नियम आधी लागू करू.
\begin{prooftree}
\AxiomC{}
\UnaryInf$ \lnot !A(a) \fCenter \lnot \lforall[x][!A(x)]$
\RightLabel{\LeftR{\lexists}}
\UnaryInf$ \lexists[x][\lnot !A(x)] \fCenter \lnot \lforall[x][!A(x)]$
\end{prooftree}
\LeftR{\lnot} आणि \RightR{\lnot} नियम उलट्या दिशेने लागू केल्यावर पुढील
निष्पत्ती मिळते:
\begin{prooftree}
\AxiomC{}
\UnaryInf$\lforall[x][!A(x)] \fCenter !A(a)$
\RightLabel{\LeftR{\lnot}}
\UnaryInf$\lnot !A(a), \lforall[x][!A(x)] \fCenter $
\RightLabel{\LeftR{\Exchange}}
\UnaryInf$\lforall[x][!A(x)], \lnot !A(a) \fCenter $
\RightLabel{\RightR{\lnot}}
\UnaryInf$ \lnot !A(a) \fCenter \lnot \lforall[x] !A(x)$
\RightLabel{\LeftR{\lexists}}
\UnaryInf$ \lexists[x] \lnot !A(x) \fCenter \lnot \lforall[x] !A(x)$
\end{prooftree}
या टप्प्यावर \LeftR{\forall} नियम लागू करणे हाच आपल्यापुढील एकमेव पर्याय
आहे. या नियमाला आयगेन चराचे बंधन लागू नसल्यामुळे आपला मार्ग मोकळा आहे.
आपल्याला आरंभीची क्रमवर्ती (जिचे स्वरूप $!A(a) \Sequent !A(a)$ असे आहे)
मिळवण्याचा प्रयत्न करायचा आहे, हे आठवा; म्हणून नियम लागू करताना $!A$ मध्ये
घालायचे पद म्हणून $a$ निवडले पाहिजे.
\begin{prooftree}
\Axiom$!A(a) \fCenter !A(a)$
\RightLabel{\LeftR{\lforall}}
\UnaryInf$\lforall[x][!A(x)] \fCenter !A(a)$
\RightLabel{\LeftR{\lnot}}
\UnaryInf$\lnot !A(a), \lforall[x][!A(x)] \fCenter $
\RightLabel{\LeftR{\Exchange}}
\UnaryInf$\lforall[x][!A(x)], \lnot !A(a) \fCenter $
\RightLabel{\RightR{\lnot}}
\UnaryInf$ \lnot !A(a) \fCenter \lnot \lforall[x][!A(x)]$
\RightLabel{\LeftR{\lexists}}
\UnaryInf$ \lexists[x][ \lnot !A(x)] \fCenter \lnot \lforall[x][!A(x)]$
\end{prooftree}
विशेषतः संख्यापक हाताळताना, या टप्प्यावर आयगेन चराच्या अटीचा भंग झालेला नाही
ना, हे पुन्हा तपासणे महत्त्वाचे असते. आयगेन चराची अट लागू असलेला आपण वापरलेला
एकमेव नियम \LeftR{\exists} होता आणि आयगेन चर~$a$ त्याच्या खालच्या क्रमवर्तीत
(म्हणजे अंतिम क्रमवर्तीत) येत नाही; म्हणून ही योग्य !!{derivation} आहे.
\end{ex}

\begin{prob}
पुढील क्रमवर्तींसाठी !!{derivation}s तयार करा:
\begin{enumerate}
\item $\Sequent (\lforall[x][!A(x)] \land \lforall[y][!B(y)]) \lif
\lforall[z][(!A(z) \land !B(z))]$.
\item $\Sequent (\lexists[x][!A(x)] \lor \lexists[y][!B(y)]) \lif
\lexists[z][(!A(z) \lor !B(z))]$.
\item $\lforall[x][(!A(x) \lif !B)] \Sequent \lexists[y][!A(y)] \lif !B$.
\item $\lforall[x][\lnot !A(x)] \Sequent \lnot\lexists[x][!A(x)]$.
\item $\Sequent \lnot\lexists[x][!A(x)] \lif \lforall[x][\lnot !A(x)]$.
\item $\Sequent \lnot\lexists[x][\lforall[y][((!A(x,y) \lif \lnot
!A(y,y)) \land (\lnot !A(y,y) \lif !A(x,y)))]]$.
\end{enumerate}
\end{prob}

\begin{prob}
पुढील क्रमवर्तींसाठी !!{derivation}s तयार करा:
\begin{enumerate}
\item $\Sequent \lnot\lforall[x][!A(x)] \lif \lexists[x][\lnot!A(x)]$.
\item $(\lforall[x][!A(x)] \lif !B) \Sequent \lexists[y][(!A(y) \lif !B)]$.
\item $\Sequent \lexists[x][(!A(x) \lif \lforall[y][!A(y)])]$.
\end{enumerate}
(या सर्वांसाठी $\RightR{\Contraction}$~नियम आवश्यक आहे.)
\end{prob}

\end{document}
